\textbf{Время:} 01:15

\textbf{Тезис:} двум неразличимым вариантам со схемы соответствуют две амплитуды, которые нужно сложить.

\textbf{На экране:} первая строка — прямой вариант a; вторая — перестановка b; внизу — сумма и обозначения.

\textbf{Текст:} Запишем формулами те же два варианта, которые были на схеме. В первой строке из A выходит пион с импульсом p-один, из B — с импульсом p-два. Каждой свободной волне соответствует комплексная экспонента. Для пары перемножаем две волны, поэтому в показателе складываются два слагаемых.

Во второй строке точки A и B те же, но импульсы поменялись местами. Это амплитуда b. В нашем упрощённом примере модули обеих амплитуд равны единице; отличаются их фазы.

Для тождественных бозонов складываем a и b со знаком плюс и делим на корень из двух. На следующем слайде этот множитель даст единичный уровень без интерференционного члена.

Внизу указаны обозначения: p объединяет энергию и импульс, x — время и положение. Произведение p на x означает энергию на ct минус обычное скалярное произведение импульса и координаты. Деление на h-bar-c делает показатель экспоненты безразмерным.

\textbf{Переход:} Амплитуду пары записали. Теперь перемножим её с комплексно сопряжённой и посмотрим, где появится интерференция.

\textbf{Справка:} это относительные плосковолновые амплитуды при близких импульсах и одинаковых модулях. Взаимодействием пионов после испускания пока пренебрегаем; источник рассматривается как хаотический. Здесь \(p^\mu=(E,\mathbf p)\) при \(c=1\), \(x^\mu=(ct,\mathbf r)\), метрика (+,−,−,−). Четырёхимпульс численно задаётся в ГэВ, координаты — в фм, \(\hc=0.19732698\) ГэВ·фм. Это соглашение не меняет экспериментальных подписей осей импульса в ГэВ/\(c\).

\textbf{Источники:} Lisa et al. (2005), §2.1–2.2 и рис. 1; собственная запись с полной разностью импульсов.
